%! Skills Focus: Selecting, Implementing, and Communicating Inference Procedures — Unit 7, Lesson 7.10 — Exit Ticket (Answer Key)
\documentclass[10pt]{article}
\usepackage{apstats-article}
\usepackage{apstats-key}   % pulls in -boxes; do NOT also load -boxes
\newcommand{\TallMath}[1]{$\displaystyle #1\rule[-1.4em]{0pt}{3.2em}$}

\begin{document}

\pageheader{Unit 7, Lesson 7.10}{Exit Ticket --- Answer Key}
\namedateperiod

\vspace{0.15in}

\textit{For each scenario: (a) name the correct inference procedure, and (b) state one key
condition that must be verified. No computation is required.}

\vspace{0.10in}
\begin{enumerate}

\item A researcher surveys 150 randomly selected registered voters and records whether each
      supports a ballot measure. She wants to \emph{estimate} the true proportion of all
      registered voters who support it.

      \begin{enumerate}[label=\alph*.]
        \item Procedure: \ans{one-sample $z$-interval for $p$}
        \item One key condition to verify: \ansline{Large Counts: $np \ge 10$ and $n(1-p) \ge 10$ (use $\hat p$ in place of $p$). Accept also: Random (random sample of 150 voters).}
      \end{enumerate}

\vspace{0.12in}

\item A teacher randomly assigns 32 students to one of two tutoring methods and records each
      student's score on a post-test. She wants to test whether the \emph{mean} post-test score
      differs between the two methods.

      \begin{enumerate}[label=\alph*.]
        \item Procedure: \ans{two-sample $t$-test for $\mu_1 - \mu_2$}
        \item One key condition to verify: \ansline{Normal/Large Samples: $n_1$ and $n_2$ are each $\ge 30$, or sample distributions show no strong skewness or outliers. (With 32 students split between two groups, $n_i$ may be $< 30$; check distribution shape.)}
      \end{enumerate}

\vspace{0.12in}

\item A confidence interval for $\mu_1 - \mu_2$ is computed to be $(1.2,\; 4.8)$.\\
      Does the interval provide evidence that $\mu_1 \ne \mu_2$? \ans{Yes}\\
      Explain: \ansline{The entire interval is above 0, so 0 is not a plausible value for $\mu_1 - \mu_2$. This is consistent with rejecting $H_0: \mu_1 - \mu_2 = 0$ at the corresponding significance level.}

\end{enumerate}

\begin{teachernote}
Item 3 connects the CI to the two-sided significance test. The key insight is that if 0 is
not in the interval, we have evidence of a real difference. Emphasize that ``the interval
does not capture 0'' is the AP-preferred phrasing.
\end{teachernote}

\end{document}
